% !TEX root = ../00_Main/Main.tex
\documentclass[../00_Main/Main.tex]{subfiles}
\begin{document}

\label{sec:soa}

El entorno de partida, PES~\citep{BCINE2022}, fue desarrollado por el
BCI-NE Lab de la Universidad de Essex para estudiar decisiones tomadas en
conjunto por personas y agentes artificiales~\citep{Nijjar2025,ds004477}.
Esta tesis utiliza \'unicamente el entorno; todas las decisiones son
tomadas por agentes artificiales.

\subsection{Transformers en aprendizaje por refuerzo}\label{sec:soa-trf}

En RL, el Transformer se ha empleado de dos maneras. El \emph{Decision
Transformer}~\citep{Chen2021} aprende de trayectorias ya registradas y
genera acciones condicionadas al retorno deseado. La segunda lo emplea
como codificador de la historia en un agente que aprende interactuando
con el entorno: Parisotto \emph{et al.}~\citep{Parisotto2020} mostraron
que en ese caso el Transformer est\'andar es dif\'icil de entrenar y que,
con normalizaci\'on de capa antes de cada subcapa y compuertas en las
conexiones residuales, alcanza o supera el desempe\~no de las LSTM en
tareas que requieren memoria. El Transformer de esta tesis sigue el segundo enfoque: es un DQN
cuyo codificador procesa los \'ultimos estados. Como el estado de mPES no
incluye la severidad de las ciudades ya atendidas ni la longitud de la
secuencia (Secci\'on~\ref{sec:env-dyn}), comparar modelos con ventana
(DQN recurrente y Transformer) y sin ella (DQN) indica si esa
informaci\'on mejora las decisiones.

\subsection{Ensambles en aprendizaje por refuerzo}

Lakshminarayanan \emph{et al.}~\citep{Lakshminarayanan2017} mostraron
que promediar redes de la misma arquitectura, entrenadas desde
inicializaciones distintas, mejora la estimaci\'on de la incertidumbre.
Wiering y van Hasselt~\citep{Wiering2008} combinaron las pol\'iticas de
cinco algoritmos distintos, entre ellos Q-learning, Sarsa y
actor--cr\'itico, mediante votaci\'on y combinaciones de distribuciones de
Boltzmann. Los ensambles de esta
tesis siguen esa l\'inea, pero combinan modelos ya entrenados y ponderan
a cada uno por la entrop\'ia de su distribuci\'on de acciones.

\subsection{Generalizaci\'on bajo variaciones controladas}

Henderson \emph{et al.}~\citep{Henderson2018} encontraron que los
algoritmos de RL profundo son sensibles a la semilla, a los
hiperpar\'ametros y a detalles de implementaci\'on; Cobbe
\emph{et al.}~\citep{Cobbe2020} mostraron que agentes con buen
rendimiento en los niveles de entrenamiento no generalizaban
a niveles nuevos producidos con la misma regla, y Kirk
\emph{et al.}~\citep{Kirk2023} recomiendan evaluar sobre familias de
variaciones controladas del entorno. Siguiendo esa recomendaci\'on, cada
modelo se eval\'ua en $21$ escenarios de generalizaci\'on, tres de
ellos por encima de las cotas del entorno
(Secci\'on~\ref{sec:scenario-catalogue}).

\biblio % Needed for referencing to working when compiling individual subfiles - Do not remove
\end{document}
